\section{交错与绝对收敛}
\label{sec:02-Calculus/05-Sequences-and-Series/04}

\[
1-\frac12+\frac13-\frac14+\frac15-\cdots=\ln 2\approx0.6931.
\]

现在把这些项重新排个顺序，改成两个正项跟一个负项的节奏：

\[
1+\frac13-\frac12+\frac15+\frac17-\frac14+\frac19+\frac1{11}-\frac16+\cdots=\frac32\ln 2\approx1.0397.
\]

没有一项被改动，没有一项被丢掉，也没有一项被重复使用——每个分数照旧只出现一次，只是排队的次序换了。和却变了。

而且 \(\tfrac32\ln2\) 没有什么特殊地位。换一种排法可以让它等于 0，等于 \(-100\)，等于圆周率，或者干脆让部分和上下振荡不收敛。这一节要解释这件事怎么发生，以及什么条件能挡住它。

\subsection{正负交替为何能救回一个发散的级数}

去掉符号，\(\sum1/n\) 发散；加上符号，它收敛。收敛的来源不是衰减，是抵消。

对 \(\sum(-1)^{n-1}b_n\)（\(b_n\ge0\)），若 \(b_n\) 最终单调递减且趋于零，则级数收敛。看部分和的走法就明白了：加一个正项往右跳，加一个负项往左跳，每次跳的距离是下一项的大小，而这个距离在缩短。于是奇数部分和从上方递减，偶数部分和从下方递增，两列一个个套进对方内部。

\begin{center}
\begin{tikzpicture}[x=0.85cm,y=6.2cm,font=\small,>=stealth]
  \draw[->,black!60] (0.5,0.45) -- (13.0,0.45);
  \node[below,black!60] at (12.95,0.443) {\(N\)};
  \draw[black!55,dashed] (0.5,0.6931) -- (12.2,0.6931);
  \node[left,black!70] at (0.45,0.6931) {\(\ln 2\)};
  \draw[black!40,thick]
    (1,1.0) -- (2,0.5) -- (3,0.8333) -- (4,0.5833) -- (5,0.7833) -- (6,0.6167)
    -- (7,0.7595) -- (8,0.6345) -- (9,0.7456) -- (10,0.6456) -- (11,0.7365)
    -- (12,0.6532);
  \fill[red!65!black] (1,1.0) circle (1.6pt);
  \fill[red!65!black] (3,0.8333) circle (1.6pt);
  \fill[red!65!black] (5,0.7833) circle (1.6pt);
  \fill[red!65!black] (7,0.7595) circle (1.6pt);
  \fill[red!65!black] (9,0.7456) circle (1.6pt);
  \fill[red!65!black] (11,0.7365) circle (1.6pt);
  \fill[blue!65!black] (2,0.5) circle (1.6pt);
  \fill[blue!65!black] (4,0.5833) circle (1.6pt);
  \fill[blue!65!black] (6,0.6167) circle (1.6pt);
  \fill[blue!65!black] (8,0.6345) circle (1.6pt);
  \fill[blue!65!black] (10,0.6456) circle (1.6pt);
  \fill[blue!65!black] (12,0.6532) circle (1.6pt);
  \draw[<->,black!75] (12.55,0.6532) -- (12.55,0.7365);
  \node[right,black!75,font=\footnotesize] at (12.6,0.695) {\(b_{12}\)};
  \node[right,red!65!black,font=\footnotesize] at (5.2,0.86) {奇数项部分和，自上而下};
  \node[right,blue!65!black,font=\footnotesize] at (5.2,0.52) {偶数项部分和，自下而上};
\end{tikzpicture}

\emph{两列部分和从两侧同时收窄，把极限夹在中间。任何时刻，真值都被关在最近的一对部分和之间，而这一对之间的距离恰好是下一项。}
\end{center}

严格论证就是把图上的话逐条兑现：奇数列单调递减有下界，偶数列单调递增有上界，各自由第一节的单调有界定理得到极限；两列之差是 \(b_{N+1}\to0\)，所以两个极限相等。三个条件缺一不可，\(b_n\) 不趋零或长期不递减时定理失效——有限个例外无所谓，从足够大的下标重新起算即可。

夹住这件事还有个直接的红利：用前 \(N\) 项近似时

\[
\abs{S-S_N}\le b_{N+1},
\]

而且误差的符号与第一个被丢掉的项相同。要把误差压到 \(10^{-3}\) 以下，看一眼下一项够不够小就行，不必另做估计。这在正项级数那边是奢望，那里得靠积分判别去围一个界。

\subsection{绝对收敛：把抵消这条路彻底关掉}

若 \(\sum\abs{a_n}\) 收敛，称 \(\sum a_n\) 绝对收敛。此时原级数必然收敛：由 Cauchy 准则，尾段的和满足 \(\bigl|\sum_{k=n+1}^{m}a_k\bigr|\le\sum_{k=n+1}^{m}\abs{a_k}\)，右边能被压到任意小，左边跟着被压住。

若 \(\sum a_n\) 收敛而 \(\sum\abs{a_n}\) 发散，称条件收敛。交错调和级数正是这一类：收敛靠的完全是符号的编排，把符号抹掉就散架。

区别可以看得更清楚一点。记

\[
a_n^+=\max(a_n,0),\qquad a_n^-=\max(-a_n,0),
\]

于是 \(a_n=a_n^+-a_n^-\) 而 \(\abs{a_n}=a_n^++a_n^-\)，两个新数列都非负。绝对收敛就是 \(\sum a_n^+\) 与 \(\sum a_n^-\) 都有限：正的总量有限，负的总量有限，各自的账都平得了。

条件收敛则必须两边都无限。理由是排除法：若两个都有限，\(\sum\abs{a_n}\) 有限，那是绝对收敛；若只有一个有限，\(\sum a_n=\sum a_n^+-\sum a_n^-\) 就是有限减无限，原级数发散。剩下唯一可能是两个都发散到无穷，级数的和只是这两个无穷之间某种精心安排的差。

到这里，开头那件怪事已经有了解释的入口：既然正的部分和负的部分各自都是取之不尽的，那么先取多少正的、再取多少负的，就成了可以操纵的自由度。

\subsection{条件收敛的级数可以被排成任何数}

\begin{theorem}[Riemann 重排定理]
若 \(\sum a_n\) 条件收敛，则对任意实数 \(L\)，存在项的一个重排使新级数收敛到 \(L\)；也存在重排使新级数发散到 \(+\infty\)、发散到 \(-\infty\)，或者部分和振荡而无极限。
\end{theorem}

构造过程是一个贪心算法，短得可以在这里走完。把非负项按原顺序抽出来记为 \(p_1,p_2,\dots\)，把负项抽出来记为 \(q_1,q_2,\dots\)，上一小节已经说明 \(\sum p_k=+\infty\) 且 \(\sum q_k=-\infty\)。

要凑到目标 \(L\)：先依次取正项，直到部分和第一次超过 \(L\) 就停手；接着依次取负项，直到部分和第一次落到 \(L\) 以下再停；然后回去继续取正项，如此往复。每一轮都能在有限步内完成，因为两边的储备各自发散，要多少有多少。每一项迟早会被用到，所以这确实是原级数的一个重排。

至于为什么收敛到 \(L\)：每次转向时，部分和越过 \(L\) 的幅度不超过刚刚放进去的那一项。而 \(a_n\to0\)（条件收敛蕴含这一点），被放进去的项越来越小，超调量随之趋于零。部分和于是围着 \(L\) 越箍越紧。把目标从固定的 \(L\) 换成一个趋于无穷的移动靶，同一套动作就给出发散的重排。

对照之下，绝对收敛的级数任意重排后和不变。原因也在正负部分上：\(\sum a_n^+\) 与 \(\sum a_n^-\) 都是有限的正项级数，而正项级数的和等于所有有限子和的上确界——上确界只看被加进来的项是哪些，不看它们的先后。两个上确界都不动，差也不动。

所以，加法的交换律并没有在无穷处失效。失效的是把``每个有限部分和都可以随意换序''直接推广成``极限也不会变''。重排改变的不是任何一个有限和，而是部分和逼近极限的整条路径。

\subsection{求和顺序无关是一个需要买的性质}

浮点求和给这条抽象结论配了一个日常版本。有限个浮点数相加不满足结合律：先把大数加完再加小数，小数会被舍入吞掉；反过来先累积小数，结果可能差出几位有效数字。多线程 reduce、分块 all-reduce、不同批大小下的梯度累加，都会因为求和顺序不同而给出不同的数值——同一份数据、同一段代码，换个并行度就复现不出上次的结果。

浮点的机制是舍入而不是重排定理，两者不能混为一谈。但它们指向同一条经验：``求和顺序不影响结果''从来不是加法自带的，它是要靠条件换来的。数学里的条件是绝对收敛，数值计算里的条件是各项量级接近或者使用补偿求和。哪一条都不满足时，顺序就是结果的一部分。相关的误差机制见 \hyperref[ch:061]{``浮点数与误差：计算机为什么不等于实数系统''}。

再往前一步，绝对收敛的地位会在测度论里被重新提出来：Lebesgue 积分定义可积时要求 \(\int\abs f<\infty\)，理由和这里一模一样——只有正部与负部各自有限，积分的值才不依赖你从哪里开始累积。\hyperref[ch:101]{``测度、积分与条件期望：为现代概率回填地基''} 会接手这条线索。

\subsection{什么时候可以放心，什么时候要盯住顺序}

判断顺序：先查 \(\sum\abs{a_n}\)。它收敛，上一节的全套正项判别法都能用上，而且此后重排、分组、与另一个绝对收敛级数逐项相乘，都不会出事。它发散而原级数收敛，那就是条件收敛，此时手上的和是与特定顺序绑定的一个数，任何改变累加次序的操作都必须重新论证。

交错级数判别法只在 \(b_n\) 递减且趋零时给结论，它给的是收敛性和误差界，不是绝对收敛。误差界 \(b_{N+1}\) 本身也别用错：它控制的是截断，不控制舍入，更不控制你把项换了顺序之后的行为。

至此讨论的都是一列固定的数。若把项写成含参数的形式，比如让第 \(n\) 项是 \(c_nx^n\)，同一个式子对不同的 \(x\) 就会给出不同的级数，有的收敛有的发散。下一节处理这种情形，并说明为什么``形式上写得出来''和``在这一点有意义''是两件必须分开的事。
